% !TEX program = xelatex
% Cards format — auto-generated from complex.tex by cards-convert.mjs
\documentclass[11pt]{article}

\usepackage{fontspec}
\usepackage{unicode-math}
% Latin Modern by OTF filename (not font name): xelatex on macOS (Core
% Text, no fontconfig) cannot resolve "Latin Modern Roman" by name, which
% silently produced ~4KB broken card PDFs. kpathsea finds the OTFs in
% texmf-dist on both macOS and the oven's Linux xelatex.
\setmainfont{lmroman10-regular.otf}[
  BoldFont=lmroman10-bold.otf,
  ItalicFont=lmroman10-italic.otf,
  BoldItalicFont=lmroman10-bolditalic.otf,
]
\setsansfont{lmsans10-regular.otf}[
  BoldFont=lmsans10-bold.otf,
  ItalicFont=lmsans10-oblique.otf,
  BoldItalicFont=lmsans10-boldoblique.otf,
]
\setmonofont{lmmono10-regular.otf}[Scale=0.88,
  BoldFont=lmmonolt10-bold.otf,
  ItalicFont=lmmono10-italic.otf,
]


\usepackage{graphicx}
\graphicspath{{figures/}{../../papers/arxiv-ac/figures/}}
\usepackage{booktabs}
\usepackage{tabularx}
\usepackage{adjustbox}
\usepackage{ragged2e}
\usepackage{microtype}
\usepackage{natbib}



\makeatletter
\def\input@path{{../}}
\makeatother
\usepackage{ac-paper-cards}


\hypersetup{
  pdftitle={Sucking on the Complex: Platform Hegemony, Critique-as-Content, and the Need for Anti-Environments},
}

\renewcommand{\acpdfbase}{sucking-on-the-complex-26-arxiv}
\begin{document}

% ============================================================
% TITLE CARD
% ============================================================
\thispagestyle{empty}
\vspace*{0.3em}
\begin{center}
\includegraphics[width=0.62\textwidth]{figures/cover}\par\vspace{0.3em}
{\acbold\fontsize{18pt}{22pt}\selectfont\color{acdark} Sucking on the Complex}\par
\vspace{0.1em}
{\fontsize{9pt}{11pt}\selectfont\color{acpink} Platform Hegemony, Critique-as-Content, and the Need for Anti-Environments}\par
\vspace{0.4em}
{\normalsize\color{cyan!70!blue}\href{https://prompt.ac/@jeffrey}{\textbf{@jeffrey}}}\par
{\small\color{acgray} Aesthetic.Computer}\par
{\small\color{acgray} ORCID: \href{https://orcid.org/0009-0007-4460-4913}{0009-0007-4460-4913}}\par
\vspace{0.4em}
\rule{0.5\textwidth}{0.5pt}\par
\vspace{0.15em}
\colorbox{yellow!60}{\small\color{red!80!black}\textbf{\textit{working draft --- not for citation}}}\par
\vspace{0.1em}
{\footnotesize\color{acgray} March 2026 · \href{https://github.com/whistlegraph/aesthetic-computer/commit/b3a48ec6a}{b3a48ec6a}}\par
\vspace{0.1em}
{\footnotesize\color{acgray}\href{https://papers.aesthetic.computer/sucking-on-the-complex-26-arxiv-da.pdf}{Dansk} · \href{https://papers.aesthetic.computer/sucking-on-the-complex-26-arxiv-es.pdf}{Español} · \href{https://papers.aesthetic.computer/sucking-on-the-complex-26-arxiv-zh.pdf}{{\accjk 中文}} · \href{https://papers.aesthetic.computer/sucking-on-the-complex-26-arxiv-ja.pdf}{{\accjk 日本語}}}\par
\end{center}

% ============================================================
% INDEX CARD
% ============================================================
\cardindex

% ============================================================
% BODY
% ============================================================
\section{The Engine and the Outside}

The standard framing of platform critique says: the complex is totalized, the artist who feeds it is complicit, and the honest move is refusal. This paper was originally drafted in that register. The present version is not. Five years of building \ac{}~\citep{scudder2026ac} alongside a decade of running @whistlegraph~\citep{scudder2026whistlegraph} have taught that the register is wrong in a specific way: it misreads what platforms are \emph{for}, from the artist's side.

Platforms are engines. They move attention at a speed and volume no institutional pipeline approaches, and they do this reliably enough that an artist can treat them as infrastructure in the way one treats electricity or a postal route. An artist who cares about building a practice that outlives the platform has a choice the critique literature almost never describes: ride the engine deliberately, using its reach to prime an audience in a formal grammar the platform did not author and cannot reproduce, and route that audience toward a practice and a \emph{userdom} that live on a different stack. This is not refusal. It is closer to extraction: using the pipe to deliver a payload the pipe cannot inspect.

Whistlegraph was such a payload. The form emerged from Radical Digital Painting~\citep{scudder2017manifesto}, 65 pre-pandemic lecture-performances including the 2018 Goodiepal \& Pals European tour~\citep{scudder2026goodiepal} and the 35th Chaos Communication Congress~\citep{scudder2018rdp35c3}, and a decade of drawing-as-singing practice that predated the account. TikTok did not make the whistlegraph viral; the whistlegraph made itself viral on TikTok, because the form's structural property---reproducibility, mark-for-syllable correspondence, the score that teaches itself~\citep{scudder2026whistlegraph}---was already there. The account's 2.6 million followers are not the work. They are the primed audience for what comes after the form's distribution phase: a \emph{userdom}---a population organized around a creative computing system they can run, fork, and perform inside, rather than an audience organized around content they can only scroll.

What follows is not a survey of platform hegemony (that diagnosis is well-rehearsed; \S2 compresses it) nor a fresh indictment of institutional critique-as-content (\S3). It is an account of what five years of building the downstream stack has produced, what structural commitments an anti-environment for creative computing rests on, how the engine primes a userdom for it, and what it costs to keep the circle open. The paper you are reading is itself one of those commitments. It is written from inside the circle, cites the rest of the circle's documentation~\citep{scudder2026ac, scudder2026kidlisp, scudder2026notepat, scudder2026os, scudder2026pieces, scudder2026urltradition, scudder2026archaeology, scudder2026whistlegraph, scudder2026goodiepal}, and is distributed through the research platter the circle has produced.

\section{The Complex}

Meta (Facebook, Instagram, Threads, WhatsApp) and ByteDance (TikTok, Douyin) together mediate the cultural visibility of nearly every working artist on earth. Instagram alone has over three billion monthly active users; it is the primary discovery platform for galleries, curators, collectors, and institutions. TikTok's short-form video format has become the default distribution mechanism for performance, studio process, and art education. Together they constitute what \citet{nieborg2018platformization} call the ``platformization of cultural production'': the contingent cultural commodity, shaped not by artistic intention but by the platform's sorting logic. For emerging and mid-career artists, platform metrics---followers, engagement rate, post frequency---have become proxy measurements for cultural relevance. The price of visibility is data: every interaction feeds the advertising machine that funds the platform~\citep{zuboff2019surveillance, srnicek2017platform}.

The algorithm does not passively display art. It actively shapes what art gets made. When an artist posts a painting and it receives high engagement, the algorithm shows it to more people; the artist, seeing the response, makes more paintings like it. This is not censorship---it is behavioral nudging, the core mechanism \citet{zuboff2019surveillance} identifies as surveillance capitalism: the prediction and modification of human behavior at scale, for profit. And on Instagram, a painting by Richter occupies the same 1080-pixel square as a photograph of someone's lunch. The platform makes no distinction between the work of decades and the work of seconds. This is not a failure of curation; it is the platform's design~\citep{chun2016updating}. Engagement is the only metric. \citet{lovink2011networks} calls the resulting structure the ``network without a cause'': a system that connects everything and contextualizes nothing. \citet{noble2018algorithms} and \citet{benjamin2019race} show that the ranking apparatus underneath the flattening is not neutral; it encodes and amplifies existing hierarchies. \citet{terranova2000free} identified the consequence twenty-five years ago as ``free labor'': unwaged cultural production that generates value for platforms. Artists are the attention economy's most dedicated free laborers.

The deepest frame for this is Attali's. In \emph{Noise}~\citep{attali1985noise}, he argues that the concert-hall orchestra is an instrument of state power: it arranges bodies into hierarchical positions under a conductor's authority, demands synchronized obedience to a written score, and produces a unified output from controlled individual labor. This is the \emph{representation} era of music---music as spectacle, as proof that power can organize noise into harmony. The tech platform is the orchestra of the twenty-first century. The founder is the conductor. The algorithm is the score. Users, creators, and artists are the musicians---each playing their part, each believing they are expressing themselves, each producing value that accrues to the institution. The platform's rhetoric of ``giving everyone a voice'' is representation-era language dressed in composition-era clothing. Everyone has a voice, but the algorithm decides who is heard. Attali glorified what was supposed to come after representation---\emph{composition}, a mode where the distinction between performer and audience dissolves, where music is made for the joy of making it. The platform is not composition. It is representation at planetary scale, with the additional extraction of behavioral data that even the court orchestra never achieved. The tech startup is structurally closer to the Hapsburg court orchestra than to the folk song. The difference is that the court paid its musicians.

The diagnosis is not original and does not need to be. What follows is.

\section{Critique-as-Content}

The dominant artistic response to platform hegemony has been to make work \emph{about} it. This work is exhibited in the world's most prestigious cultural institutions. It is celebrated, collected, and discussed. The structural conditions it describes remain unchanged.

The mode takes several forms. Paglen photographs classified military installations and documents the physical infrastructure of surveillance, exhibiting at the Whitney, MoMA, the Smithsonian, and the Venice Biennale; the work advertises the complex's reach without routing its audience toward anything else. \citet{steyerl2017duty} describes a world of ``duty-free art'' circulating in free-trade zones, and her video installations at the Serpentine, the Venice Biennale, and the Park Avenue Armory are analyses of the same circulation---shared on Instagram, discussed on TikTok, the critique distributed by the machine it describes (\citet{steyerl2009poor} argued that the ``poor image'' gains political power through its low-resolution circulation; on the feed, every image generates the same behavioral data). \citet{blas2014informatic} developed the ``Facial Weaponization Suite'' (2011--2014), collective masks that defeat facial recognition, shown at the Whitechapel and Tate Modern---but the mask is an art object in a vitrine, not a tool you can wear. Mindy Seu's \emph{Cyberfeminism Index}~\citep{seu2023cyberfeminism} catalogs over 700 resources tracing cyberfeminist resistance, much of it distributed through Instagram Stories: a feminist archive of digital resistance made legible to the art world through Meta's surveillance infrastructure.

\citet{ulman2014excellences} is the canonical case. Her five-month Instagram performance in 2014---a scripted manipulation of the platform's logic of aspiration and self-display, later exhibited at Tate Modern in 2016---succeeded precisely because it was indistinguishable from content. The platform did not know it was being used as material. The behavioral data it generated was identical to that of genuine lifestyle posts. This is the structural point, not an indictment of any individual artist: the algorithm does not distinguish critique from endorsement. Paglen's photograph of an NSA facility and a landscape painting both generate engagement; both feed the machine. The pipeline is: make content for the feed $\rightarrow$ get noticed by curators (who found you on Instagram) $\rightarrow$ exhibit the same content or a critique of it in institutions whose audiences were also found on Instagram. \citet{fraser2005critique} identified this loop in institutional critique three decades ago; platform hegemony is the same structural problem with surveillance capitalism layered on top.

The specific failure of critique-as-content is not that it operates inside the platform---everything does---but that it \emph{critiques without cultivating}. It uses the platform's reach to amplify a diagnosis and stops there. It does not use the reach to prime an audience for a practice that lives on a different stack. The platform is perfectly willing to distribute the diagnosis; the diagnosis is content, and content is what the platform distributes. What the platform cannot distribute is a userdom. That is work that has to be done deliberately, with a form the platform can carry but not reproduce.

\section{Free Spaces of Play}

Art does not grow in feeds. \citet{winnicott1971playing} called the space where creative experience becomes possible \emph{potential space}---a transitional area of experience between inner and outer reality, fragile, requiring safety, continuity, and freedom from intrusion. \citet{huizinga1938homo} described the ``magic circle'' of play: a bounded space with its own rules, separated from ordinary life. Inside the magic circle, the rules of the outside world are suspended. A studio, when it is working, is a magic circle. A rehearsal room is a magic circle. A sketchbook is a magic circle.

The platform has no magic circle. Every creative act performed on Instagram is immediately measured, sorted, ranked, and fed back as engagement data. The potential space collapses into metrics the moment the work is posted. \citet{fisher2009capitalist} called the pervasive sense that there is no alternative \emph{capitalist realism}: the realism is not in the content but in the inability to imagine otherwise. Art school teaches the shape of a magic circle; platforms do not. This paper is an argument for imagining otherwise---and a report on what otherwise looks like once it is built.

\section{Anti-Environments}

McLuhan argued that you cannot perceive an environment from inside it~\citep{mcluhan1964understanding, mcluhan1969playboy}. An environment is invisible to its inhabitants precisely because it is total. The only way to make an environment visible is to construct an \emph{anti-environment}: a counter-structure whose difference from the dominant environment reveals the environment's shape---what \citet{shklovsky1917art} called \emph{defamiliarization}, the capacity of art to make the habitual strange.

This paper uses \emph{anti-environment} in its structural sense: not a protest, not a showroom, not a refusal. A working alternative. An instrument. A bounded, rule-carrying structure that is built alongside the dominant environment and routes creative labor through its own stack rather than through the platform's. The relationship to the platform is neither oppositional nor abstentionist. An anti-environment can use a platform's reach without inheriting its logic, in the way that mail art used the postal service without becoming a post office. Community radio, mail art networks, zine culture, Fluxus distribution circuits, and Nelson's \emph{Computer Lib}~\citep{nelson1974computerlib} all operated as anti-environments to the broadcast media of their time; \citet{lanier2018ten} arrives at a related conclusion from inside the industry. Some survived; many were absorbed. Survival depended less on refusal than on whether the anti-environment had a practice the dominant environment could not reproduce.

An anti-environment for creative computing in 2026 is a stack: a runtime, a language, an instrument, a hardware substrate, a community protocol. Each layer is a specific structural commitment that refuses platform logic at the level of the layer, not at the level of the message. The next section describes the commitments \ac{} has made, each documented in depth by a companion paper in the accompanying research platter.

\section{What Has Been Built}

Five years of operation have produced a set of structural commitments. Each is a specific refusal of platform logic at the level of the stack, and each is documented in depth by a sibling paper in the accompanying research platter. This section is the synthesis: these are not separate decisions. They are the same move.

\subsection{URL-first: the address is the interface}

Every piece on \ac{} is a URL. There is no feed, no menu, no file picker, no project list, no algorithmic surface. The prompt \emph{is} the address bar: you type a word and you are in the piece~\citep{scudder2026urltradition}. This is the HyperCard tradition re-grounded in the web, descending through Glitch's instant remix links and p5.js's Web Editor URLs. A platform has a home page you arrive at; \ac{} has a hundred addresses you type toward. The URL is not a feature. It is the medium property of the whole system.

\subsection{Piece-not-post: the bounded unit}

The unit of cultural production on \ac{} is a \emph{piece}: a self-contained interactive program~\citep{scudder2026pieces}, addressable by name, persistent across time, not algorithmically sorted. A piece is to a post what a song is to a scroll. You do not discover a piece because an algorithm surfaced it; you arrive at a piece because someone named it to you, or because you typed its address, or because you found it in the publicly searchable list. Winnicott's \emph{potential space}~\citep{winnicott1971playing} is restored at the level of the unit: each piece is bounded, ruled, and uninterrupted by engagement metrics.

\subsection{Instrument-not-profile: the interface is played}

\np{}~\citep{scudder2026notepat} is a chromatic keyboard synthesizer that runs in the browser and on bare metal. It is an instrument you play. It is not a profile you curate, not a feed you scroll, not a brand you maintain. The difference between instrument and profile is the difference between use and display. A profile is optimized to be seen; an instrument is optimized to be used. \ac{} treats the creative interface as an instrument at every level of the stack---from \np{}'s keyboard to KidLisp's REPL to the prompt itself---and this choice downgrades legibility-to-others as a design goal in favor of immediacy-to-self.

\subsection{KidLisp: a composition-era language}

KidLisp~\citep{scudder2026kidlisp} is a minimal Lisp dialect for generative art. It is a composition-era language in Attali's sense~\citep{attali1985noise}: the distinction between performer and audience dissolves inside it. The same code you write to paint a picture is the code others read to understand the picture. The same REPL you use to sketch is the REPL others use to fork. This is not a general property of Lisp---most Lisps are architect-era, built for specifying systems---but a property of KidLisp specifically, which trades expressive power for immediacy, forkability, and readability at 8-bit resolution.

\subsection{Sovereign node: the hardware is yours}

\ac{} ships as a bare-metal operating system that boots surplus hardware directly into the creative runtime~\citep{scudder2026os, scudder2026identity}. There is no cloud account, no vendor lock, no update-expiration timer, no behavioral telemetry. The laptop is a sovereign node on the network: it runs its own server, holds its own data, federates with other nodes on its own terms. This is the anti-Chromebook~\citep{scudder2026openschools}: the same \$50 hardware, flashed with software the user owns, on a protocol stack they can audit. With roughly 240 million Windows 10 end-of-life machines heading to landfill in the 2025--2026 cycle, the sovereign-node commitment is not a niche political gesture; it is a material program for what to do with the hardware.

\subsection{Pals: affinity, not followership}

The \ac{} community is organized on the Pals model, inherited from Goodiepal's tours and documented at depth in the Goodiepal paper~\citep{scudder2026goodiepal}. A Pal is not a follower. A Pal is someone who makes with you, travels with you, shares the instrument with you. The relationship is reciprocal, affinity-based, and face-to-face where possible; it is not a one-way broadcast channel. Platforms optimize for parasocial one-to-many flow because that is what advertising monetizes. Pals optimize for one-to-one depth because that is what practice rewards. The numbers are smaller. The bonds are denser. The model scales by fork, not by reach. Lewis's ``making kin with the machines''~\citep{lewis2018making}, \citet{escobar2018designs}'s designs for the pluriverse, and \citet{costanzachock2020design}'s design justice articulate the same commitment in different registers: relational rather than transactional, care rather than extraction.

\subsection{The 94-project lineage}

None of the above is new. The \ac{} repository archaeology traces the platform's technical genealogy through 94 predecessor projects spanning 2007--2020~\citep{scudder2026archaeology}---drawing tools, performance tools, development infrastructure, games, instruments. The commitments above did not arrive as a design philosophy. They accumulated as the resolution of problems that kept recurring across 94 attempts. The anti-environment is not an idea the author had. It is the shape that fell out of practicing the same practice for two decades inside media environments that were, each in their turn, total.

\section{The Engine and the Userdom}

The objection to the argument so far writes itself: if anti-environments are built outside the platform, what to make of the @whistlegraph TikTok account, which reached 2.6 million followers during the COVID-19 pandemic~\citep{scudder2026whistlegraph} and still moves roughly forty-eight thousand video views a day more than two years after the trio separated in November 2023?\footnote{TikTok Analytics snapshot of @whistlegraph, 24 April 2025 through 23 April 2026---the second full year after the trio's separation, during which no new whistlegraphs were produced by the original creators: $\sim$17.69\,million video views (median 43{,}047/day), $\sim$465{,}000 likes, $\sim$16{,}000 shares, $\sim$191{,}000 profile views. Peak single-day views in this window (208{,}000 on 27 September 2025; 198{,}000 on 16 February 2026) reflect continued organic circulation of the original archive.} Is that not evidence that the outside exists only by permission of the inside?

It is evidence of something more precise. Platforms are distribution engines. They move attention at a speed and scale no institutional pipeline approaches, and reliably enough that an artist can treat them as infrastructure. The question the standard platform-critique literature does not ask---because it is committed to an oppositional frame---is what the engine is being used to \emph{carry}. The @whistlegraph account carried a specific formal grammar: graphic notation (marks that correspond to sung syllables), interface rhetoric (performance as lecture as instrument), and a reproducibility logic (the score that teaches itself~\citep{scudder2026whistlegraph}). That grammar was developed pre-platform through Radical Digital Painting~\citep{scudder2017manifesto, scudder2018rdp35c3}, 65 lecture-performances including a 2018 European tour with Goodiepal~\citep{scudder2026goodiepal}, and a decade of drawing-as-singing practice the platform had no hand in authoring. TikTok did not make the whistlegraph; it distributed one.

What the distribution primed is a \emph{userdom}: a population of roughly 2.6 million people who learned, through watching and reproducing whistlegraphs, to recognize a specific kind of formal move---marks as instructions, interfaces as instruments, performance as pedagogy. The userdom is not the follower count. It is the primed audience available downstream, at whatever rate it can be converted into practice on a system that carries the same grammar deeper than the platform can host. \ac{} is that system~\citep{scudder2026ac}. The userdom is its seed population. The network audit~\citep{scudder2026audit} documents the current shape of the converted userdom---2{,}800+ registered handles, 16{,}000+ KidLisp programs, 93{,}000+ boot sessions---a small fraction of the primed audience and more than any comparable creative computing tool has produced without an institutional parent.

The paradox the earlier draft of this paper could not resolve is this: the largest anti-environment in the \ac{} platter was built \emph{through} the complex. The resolution is that the complex was never the adversary. It was the engine. The mistake the critique-as-content literature makes is to read every use of the platform as capitulation. The artist's actual move---the one available to anyone with a form the platform's algorithm will amplify---is to use the engine deliberately, delivering a payload the platform cannot inspect, toward a userdom the platform cannot hold. Scale is rented from the engine. Practice is owned by the userdom.

\section{Funding the Engine and the Circle}

Anti-environments are hard to fund for a structural reason: the platform economy prices attention, and the anti-environment is built to refuse attention's measurability. A piece has no engagement rate. A Pal has no cost-per-mille. An instrument has no retention curve. The sustainability paper in this series~\citep{scudder2026sustainability} surveys how 28 comparable creative computing tools have been financed---academic subsidy, corporate patronage, personal sacrifice, community crowdfunding, windfall, or no model at all---and finds that the tools most widely used are not the tools best funded. The median gap between a tool's creation and its first sustainable funding is eight years. Most tools close before crossing it.

The engine-and-userdom frame does not solve the funding problem but clarifies it. The engine is free to use and free to leave. The userdom is the asset the engine produces. The practice is where the asset is spent. The funding question is not whether to monetize the engine---the engine monetizes artists on its own terms, selling their attention to advertisers---but how to convert the userdom into a self-sustaining practice economy before the engine's next algorithm update rotates the attention elsewhere. Illich's \emph{tools for conviviality}~\citep{illich1973tools} name the destination: tools built at human scale, by hand, for use rather than consumption. The economy of such tools is still being written. This paper is one entry in its documentation.

\section{Coda}

Anti-environments are possible. They are not rare; they are simply not legible to the platform, because the platform was not built to read what they produce. This paper has described the shape of one, from inside, with the specific structural commitments it rests on written down so that others can recognize them, fork them, adapt them, or refuse them:

\begin{itemize}
  \item URL-first, not feed-first.
  \item Piece-not-post: bounded units, not scrolling streams.
  \item Instrument-not-profile: interfaces to play, not to curate.
  \item Composition-era language: code read the way it is written.
  \item Sovereign node: hardware and data you own.
  \item Affinity, not followership: scale by fork, not by reach.
  \item Engine, not adversary: use the platform's reach to prime a userdom whose practice lives elsewhere.
\end{itemize}

These are the rules of one circle. Other people can draw their own. The algorithm cannot read any of them.

\vspace{0.5em}
\noindent\textbf{ORCID:} \href{https://orcid.org/0009-0007-4460-4913}{0009-0007-4460-4913}

\bibliographystyle{plainnat}
\bibliography{references}

\end{document}
